% Part: history
% Chapter: biographies
% Section: bertrand-russell

\documentclass[../../../include/open-logic-section]{subfiles}

\begin{document}

\olfileid{his}{bio}{rus}

\olsection{برتراند راسل}

\olphoto{russell-bertrand}{Bertrand Russell}

برتراند راسل را یکی از بنیان‌گذاران فلسفهٔ تحلیلی نوین می‌دانند. راسل
که در 18 مه 1872 به دنیا آمد، نه‌تنها برای کارش در فلسفه و منطق
شناخته می‌شد، بلکه کتاب‌های عامه‌فهم بسیاری در موضوعات گوناگون نوشت.
او همچنین در سراسر زندگی کنشگری سیاسی و پرشور بود.

راسل در ترلکِ مونموت‌شر در ولز به دنیا آمد. والدینش از اشراف بریتانیا
بودند. آن‌ها آزاداندیش بودند و حتی با تندروهای آن زمان در بوستون
دوستی داشتند. متأسفانه والدین راسل در کودکی او درگذشتند و راسل را برای
زندگی نزد پدربزرگ و مادربزرگش فرستادند. آنجا تربیتی دینی یافت، درست
همان چیزی که والدینش می‌خواستند به هر قیمتی از آن جلوگیری کنند.
مادربزرگش در همهٔ امور اخلاقی بسیار سخت‌گیر بود. راسل در نوجوانی بیشتر
در خانه و نزد آموزگاران خصوصی تحصیل کرد.

تأثیر راسل بر فلسفهٔ تحلیلی، و به‌ویژه بر منطق، بسیار بزرگ است. او در
کالج ترینیتی کمبریج ریاضیات و فلسفه خواند و آنجا از آلفرد نورث وایتهدِ
ریاضی‌دان و فیلسوف تأثیر پذیرفت. راسل و وایتهد در سال 1910 نخستین جلد
\emph{Principia Mathematica} را منتشر کردند و در آن از این دیدگاه دفاع
کردند که ریاضیات فروکاستنی به منطق است. راسل سپس صدها کتاب، جستار و
جزوهٔ سیاسی منتشر کرد. او در سال 1950 جایزهٔ نوبل ادبیات را دریافت کرد.

راسل عمیقاً در سیاست و کنشگری اجتماعی درگیر بود. در جنگ جهانی اول به
سبب فعالیت‌ها و اعتراض‌های صلح‌طلبانه بازداشت شد و شش ماه به زندان
افتاد. در زندان توانست بنویسد و بخواند و ادعا کرد این تجربه را «بسیار
خوشایند» یافته است. در سراسر عمر صلح‌طلب ماند و در سال 1961 بار دیگر
به سبب حضور در گردهمایی خلع سلاح هسته‌ای زندانی شد. او در 1948 از
سانحهٔ سقوط هواپیما نیز جان سالم به در برد؛ تنها بازماندگان کسانی
بودند که در بخش سیگارکشیدن نشسته بودند. از همین رو راسل ادعا کرد
زندگی‌اش را مدیون سیگار است. راسل چهار بار ازدواج کرد، اما به روابط
خارج از ازدواج شهرت داشت. او در 2 فوریهٔ 1970 در 97سالگی در
پنریندایدرایثِ ولز درگذشت.

\begin{reading}
راسل زندگی‌نامهٔ خودنوشتش را در سه بخش نوشت که زندگی او از 1872 تا
1967 را در بر می‌گیرد \citep{Russell1967,Russell1968,Russell1969}.
مرکز پژوهشی برتراند راسل در دانشگاه مک‌مستر میزبان بایگانی‌های برتراند
راسل است. برای اطلاعات دربارهٔ جلدهای مجموعه‌آثار او، شامل نمایه‌های
قابل جست‌وجو، و پروژه‌های بایگانی، وبگاه آن‌ها را در
\citet{Duncan2015} ببینید. مقالهٔ راسل با عنوان \emph{On Denoting}
\citep{Russell1905} اثری کلاسیک در فلسفهٔ تحلیلی قرن بیستم است.

مدخل دانشنامهٔ فلسفهٔ استنفورد دربارهٔ راسل \citep{Irvine2015}
قطعه‌های صوتی از سخنان راسل دربارهٔ میل و نظریهٔ سیاسی دارد. گفت‌وگوهای
ویدئویی بسیاری با راسل در اینترنت در دسترس است. برای نمونه، برای دیدن
سخنان او دربارهٔ سیگارکشیدن و سانحهٔ هواپیما، \citet{RussellND} را
ببینید. برخی آثار راسل، از جمله
\emph{Introduction to Mathematical Philosophy}، به‌صورت کتاب صوتی
رایگان در \citet{LibriVoxND} در دسترس‌اند.
\end{reading}

\end{document}
